%% Question bank chapter file - generated by bank_tools.py build (do not edit while the bank is still being built:
%% it is regenerated from the per-test data in Question Bank\data). Every question is one bankq block:
%%   \begin{bankq}[id=..., source={...}, paper=.., qno=.., type=..., subtopic={...}, answer={...}, answer_src=key|solved,
%%                 working={...}, difficulty=1-10, diff_parts={K S I C F}, diff_note={...},
%%                 group=<passage id or empty>, also_in={duplicate copies in other tests}]
%% difficulty = min(10, K+S+I+C+F-1): K concepts 1-3, S steps 1-4, I insight 0-2, C computation 0-2, F format 0-1
%%   ...question LaTeX, options, figures (\bfig / \bfigi -> img/)...
%%   \end{bankq}
%% Passages shared by several questions are bankpassage blocks; their questions carry group=<passage id>.
\documentclass[10pt]{article}
\usepackage[a4paper,margin=1.8cm]{geometry}
\usepackage[T1]{fontenc}
\usepackage{amsmath,amssymb,mathtools}
\usepackage{graphicx}
\usepackage[export]{adjustbox}
\usepackage{xcolor,needspace,keyval}
\usepackage{enumitem}
\usepackage{booktabs,array,tabularx}
\usepackage[hidelinks]{hyperref}
\setlength{\parindent}{0pt}
\setlength{\parskip}{3pt}
\makeatletter
\define@key{bq}{id}{\def\bq@id{#1}}\define@key{bq}{source}{\def\bq@src{#1}}\define@key{bq}{paper}{\def\bq@paper{#1}}
\define@key{bq}{qno}{\def\bq@qno{#1}}\define@key{bq}{type}{\def\bq@type{#1}}\define@key{bq}{subtopic}{\def\bq@sub{#1}}
\define@key{bq}{answer}{\def\bq@ans{#1}}\define@key{bq}{answer_src}{\def\bq@asrc{#1}}\define@key{bq}{working}{\def\bq@work{#1}}
\define@key{bq}{group}{\def\bq@group{#1}}\define@key{bq}{also_in}{\def\bq@also{#1}}
\define@key{bq}{difficulty}{\def\bq@diff{#1}}\define@key{bq}{diff_parts}{\def\bq@dparts{#1}}\define@key{bq}{diff_note}{\def\bq@dnote{#1}}
\newenvironment{bankq}[1][]{\def\bq@work{}\def\bq@also{}\def\bq@group{}\def\bq@diff{?}\def\bq@dparts{}\def\bq@dnote{}\setkeys{bq}{#1}%
  \par\needspace{4\baselineskip}\medskip\noindent{\small\textbf{\bq@id}\hfill\color{gray}[\bq@src, Paper \bq@paper, Q\bq@qno]
  \bq@type\ $\cdot$ \bq@sub\ $\cdot$ \textbf{D\bq@diff}}\par\nopagebreak\relax}
  {\par{\small\color{teal}Answer: \textbf{\bq@ans}\ifx\bq@work\@empty\else\ (solved: \bq@work)\fi
  \ifx\bq@also\@empty\else\par\color{gray}Also in: \bq@also\fi}\par}
\define@key{bp}{id}{\def\bp@id{#1}}\define@key{bp}{source}{\def\bp@src{#1}}\define@key{bp}{questions}{\def\bp@qs{#1}}
\newenvironment{bankpassage}[1][]{\setkeys{bp}{#1}\par\needspace{4\baselineskip}\bigskip\noindent
  {\small\textbf{Passage \bp@id}\hfill\color{gray}[\bp@src; questions \bp@qs]}\par\nopagebreak\relax\itshape}{\par}
\makeatother
\newcommand{\bfig}[2]{\par\smallskip{\centering\includegraphics[width=#1pt,max width=0.62\linewidth,max height=0.32\textheight]{img/#2}\par}\smallskip}
\newcommand{\bfigi}[2]{\includegraphics[width=#1pt,max width=\linewidth,max height=0.32\textheight,valign=c]{img/#2}}
\newcommand{\bmissing}[1]{\textit{[figure missing in the source conversion: \detokenize{#1}]}}
\newenvironment{qoptions}{\begin{itemize}[leftmargin=1.8em,itemsep=1pt,topsep=2pt,label={}]}{\end{itemize}}
\newcommand{\opt}[2]{\item[(#1)] #2}
\providecommand{\fwstart}{}\providecommand{\fwend}{}\providecommand{\qfullsplit}{}
\newcolumntype{L}[1]{>{\raggedright\arraybackslash}p{#1}}
\newcolumntype{C}[1]{>{\centering\arraybackslash}p{#1}}
\newcolumntype{R}[1]{>{\raggedleft\arraybackslash}p{#1}}
\providecommand{\sectionhdr}[1]{}\providecommand{\subsectionhdr}[1]{}


\begin{document}
{\Large\bfseries Polymers --- Question Bank}\par\smallskip
{\small Chemistry chapter C27. Questions from all converted test papers, grouped by source series. Answers: printed key where it matches the paper, otherwise solved (answer\_src=solved, with working). Exact repeats across papers appear once, with also\_in listing the other copies.}\par\medskip\hrule

\section*{Score-I papers}
\begin{bankq}[id=SI12-P2-CHEM-S3-Q07, source={Score-I - Test 12 (09 Mar 2022)}, paper=2, qno=7, type=matching, subtopic={Examples}, answer={P$\to$2, Q$\to$1, R$\to$3}, answer_src=key, difficulty=2, diff_parts={K1 S1 I0 C0 F1}, diff_note={concepts: uses of polymers | steps: uses of nylon-6 and bakelite; use of HDPE}]
\begin{tabular}[t]{|p{0.65cm}|p{1.6cm}|p{0.65cm}|p{3.8cm}|}
\hline
 & \centering List-I & & \centering List-II\tabularnewline
\hline
(P) & Nylon-6 & (1) & Computer disc\\
(Q) & Bakelite & (2) & Tyre cords, Ropes\\
(R) & HDPE & (3) & Buckets, Dustbins \& Bottles, Pipes\\
\hline
\end{tabular}
\end{bankq}
\begin{bankq}[id=SI19-P1-CHEM-S3-Q04, source={Score-I - Test 19 (20 Nov 2022)}, paper=1, qno=4, type=numerical, subtopic={Types}, answer={6}, answer_src=key, difficulty=5, diff_parts={K1 S2 I1 C1 F1}, diff_note={concepts: count addition polymers | steps: polyethene Teflon PVC PMMA polystyrene natural rubber are addition polymers; others condensation; count}]
How many of the following are addition polymers ?\\
Polyethene, Teflon, PVC, Nylon, Bakelite, Melamine-formaldehyde resin, PMMA, Polystyrene, Natural rubber, Urea-formaldehyde resin.
\end{bankq}
\begin{bankq}[id=SI31-P2-CHEM-S1-Q17, source={Score-I - Test 31 (29 Dec 2022)}, paper=2, qno=17, type=single, subtopic={Types}, answer={C}, answer_src=key, difficulty=2, diff_parts={K1 S1 I1 C0 F0}, diff_note={concepts: name of the polymer from its repeat unit | steps: butadiene and acrylonitrile units give Buna-N}]
$\left[\mathrm{CH_2{-}CH{=}CH{-}CH_2{-}CH_2{-}}\ \shortstack{CN\\$|$\\CH}\ \right]_n$

The name of polymer is :
\begin{qoptions}
\opt{A}{Polystyrene}
\opt{B}{BuNa--S}
\opt{C}{BuNa--N}
\opt{D}{BuNa--CN}
\end{qoptions}
\end{bankq}

\section*{Score-II papers}
\begin{bankq}[id=SII01-P1-CHEM-S3-Q07, source={Score-II - Test 01 (21 Mar 2021)}, paper=1, qno=7, type=numerical, subtopic={Examples}, answer={6.00}, answer_src=key, difficulty=5, diff_parts={K1 S2 I1 C1 F1}, diff_note={concepts: count correct statements about penicillin | steps: structure features; tests; sulphur detection; count}]
How many of the following is/are correct for Penicillin.
(a) It gives Black ppt with lead Acetate. [Lassaigne Test]
(b) It gives effervescence of CO$_2$ with NaHCO$_3$
(c) Presence of Lactum in its structure.
(d) It gives Red colour in Victor Mayer Test.
(e) It forms pungent smelling compound with CHCl$_3$ + NaOH
(f) It do not give violet colour with neutral FeCl$_3$
(g) It give +ve Mustard oil test.
(h) It give Prussian blue colour in Lassaigne Test
(i) 1 mol of Penicillin releases 1 mol of H$_2$(g) on treatment with Na.
\end{bankq}
\begin{bankq}[id=SII07-P1-CHEM-S2-Q03, source={Score-II - Test 07 (10 Jul 2022)}, paper=1, qno=3, type=numerical, subtopic={Examples}, answer={5}, answer_src=key, difficulty=3, diff_parts={K1 S2 I0 C0 F1}, diff_note={concepts: condensation polymers | steps: terylene nylon 6,6 bakelite novolac are condensation polymers; cellulose is a natural addition of glucose by condensation but not counted; PVC SBR PTFE neoprene are addition; count}]
How many of the following are prepared by condensation polymerization ?\\
Dacron, Nylon-6, Nylon-66, PVC, SBR, PTFE, Bakelite, Novolac, Neoprene
\end{bankq}
\begin{bankq}[id=SII08-P2-CHEM-S2-Q10, source={Score-II - Test 08 (05 Feb 2023)}, paper=2, qno=10, type=multi, subtopic={Examples}, answer={A,B,C}, answer_src=key, difficulty=3, diff_parts={K1 S2 I0 C0 F1}, diff_note={concepts: polymers and their preparation | steps: neoprene from chloroprene; Teflon by persulphate at high pressure; thermoplastic PVC; ethene polymerisation at low density conditions}]
Among the following, the correct statement(s) about polymers is(are)
\begin{qoptions}
\opt{A}{The polymerization of chloroprene gives neoprene as polymer.}
\opt{B}{Teflon is prepared from tetrafluoroethene by heating it with persulphate catalyst at high pressures.}
\opt{C}{PVC are thermoplastic polymers.}
\opt{D}{Ethene at 350--570 K temperature and 1000--2000 atm pressure in the presence of a peroxide initiator yields high density polythene.}
\end{qoptions}
\end{bankq}
\begin{bankq}[id=SII09-P2-CHEM-S2-Q09, source={Score-II - Test 09 (12 Feb 2023)}, paper=2, qno=9, type=multi, subtopic={Examples}, answer={A,D}, answer_src=key, difficulty=6, diff_parts={K1 S3 I1 C1 F1}, diff_note={concepts: reaction sequence from the styrene-type monomer of Buna-S | steps: heavier monomer styrene; oxidation to benzoic acid; acid chloride; reaction with dimethylcadmium gives acetophenone; Baeyer-Villiger oxidation gives phenyl acetate; check each statement}]
P (The heavier monomer of Buna--S) $\xrightarrow[\text{(ii)}H_3O^+]{\text{(i)}KMnO_4,OH^-}Q\xrightarrow[CHCl_3]{SOCl_2}R\xrightarrow{(CH_3)_2Cd}S\xrightarrow{MCPBA}T$
\bfig{160.4}{SII09_p2_chem_s2_q09_1.png}
Based on the above sequence of chemical reaction which is/are CORRECT?
\begin{qoptions}
\opt{A}{P has DU = 5}
\opt{B}{S is \bfigi{47.7}{SII09_p2_chem_s2_q09_2.png}}
\opt{C}{T is \bfigi{55.8}{SII09_p2_chem_s2_q09_3.png}}
\opt{D}{S can undergo aldol condensation}
\end{qoptions}
\end{bankq}
\begin{bankq}[id=SII15-P1-CHEM-S1-Q12, source={Score-II - Test 15 (09 Feb 2024)}, paper=1, qno=12, type=single, subtopic={Examples}, answer={C}, answer_src=key, difficulty=1, diff_parts={K1 S1 I0 C0 F0}, diff_note={concepts: bakelite formation | steps: phenol with formaldehyde}]
Thermosetting polymer, Bakelite is formed by the reaction of phenol with :-
\begin{qoptions}
\opt{A}{CH$_3$CH$_2$CHO}
\opt{B}{CH$_3$CHO}
\opt{C}{HCHO}
\opt{D}{HCOOH}
\end{qoptions}
\end{bankq}
\begin{bankq}[id=SII15-P1-CHEM-S1-Q14, source={Score-II - Test 15 (09 Feb 2024)}, paper=1, qno=14, type=single, subtopic={Examples}, answer={C}, answer_src=key, difficulty=3, diff_parts={K1 S2 I1 C0 F0}, diff_note={concepts: match compounds with their tests | steps: Molisch for carbohydrates; FeCl3 violet for phenol; other tests; match}]
\par\nobreak
\begin{center}\small
\begin{tabular}{|c|p{2.6cm}|c|p{2.5cm}|}\hline
 & \multicolumn{1}{c|}{\textbf{List-I}} & & \multicolumn{1}{c|}{\textbf{List-II}}\\\hline
(P) & \bfigi{56.1}{SII15_p1_chem_s1_q14_1.png} & (1) & Product give molish test\\\hline
(Q) & \bfigi{55.9}{SII15_p1_chem_s1_q14_2.png} & (2) & Product give violet colour with neutral FeCl$_3$\\\hline
(R) & \bfigi{88.5}{SII15_p1_chem_s1_q14_3.png}\newline Starch product & (3) & L.S. of product give blue colour solution with FeSO$_4$ solution\\\hline
(S) & \bfigi{85.7}{SII15_p1_chem_s1_q14_4.png} & (4) & Product give Biurate test\\\hline
\end{tabular}
\end{center}
\begin{qoptions}
\opt{A}{P $\to$ 3;Q $\to$ 2,3;R $\to$ 1;S $\to$ 4}
\opt{B}{P $\to$ 2;Q $\to$ 2,3;R $\to$ 1;S $\to$ 3,4}
\opt{C}{P $\to$ 3;Q $\to$ 2,3;R $\to$ 1;S $\to$ 3,4}
\opt{D}{P $\to$ 3;Q $\to$ 2;R $\to$ 4;S $\to$ 1,3}
\end{qoptions}
\end{bankq}
\begin{bankq}[id=SII15-P1-CHEM-S2-Q06, source={Score-II - Test 15 (09 Feb 2024)}, paper=1, qno=6, type=numerical, subtopic={Examples}, answer={Bonus}, answer_src=solved, working={Printed key omits or corrupts this answer and the question was not solved; banked as Bonus.}, difficulty=5, diff_parts={K1 S3 I0 C1 F1}, diff_note={concepts: count correct property matches for oxides and halides | steps: ClO2 odd electron but does not dimerise; XeF8 2- exists as square antiprism; Cl2O bond angle; other items; count}]
How many of the following have been correctly matched with their property.\\
(a) ClO$_2$ $-$ Contains odd electron but does not undergo dimerization.\\
(b) XeF$_8^{-2}$ $-$ Does not exist.\\
(c) Cl$_2$O $-$ Bond angle greater than 109$^\circ$28.\\
(d) B$_2$ $-$ Diamagnetic if s-p mixing in not considered.\\
(e) NO $-$ Diamagnetic gas\\
(f) N$_2$H$_4$ - N-N bond length greater than N$_2$F$_4$.\\
(g) H$_2$S$_2$O$_8$ $-$ Contains peroxy linkage.\\
(h) XeF$_4$ $-$ d$_{x^2-y^2}$ orbital is involved in hybridisation of central atom.
\end{bankq}
\begin{bankq}[id=SII15-P2-CHEM-S3-Q07, source={Score-II - Test 15 (09 Feb 2024)}, paper=2, qno=7, type=numerical, subtopic={Examples}, answer={3}, answer_src=key, difficulty=5, diff_parts={K1 S2 I1 C1 F1}, diff_note={concepts: count condensation polymers | steps: nylon 6 6 10 and others from condensation; addition polymers excluded; count}]
Among the following number of condensation polymers are ?\\
Nylon-6, Buna-N, Buna-S, Nylon-6,6, Nylon-6,10, PVC, Polystyrene, Teflon.
\end{bankq}
\begin{bankq}[id=SII21-P2-CHEM-S2-Q10, source={Score-II - Test 21 (14 Mar 2024)}, paper=2, qno=10, type=multi, subtopic={Examples}, answer={B,C}, answer_src=key, difficulty=3, diff_parts={K1 S2 I0 C0 F1}, diff_note={concepts: statements about polymers | steps: chloroprene gives neoprene not natural rubber; Teflon from persulphate catalyst; glyptal is a condensation polymer; low density polythene is from 350-570 K high pressure}]
Among the following, the correct statement(s) about polymers is(are)
\begin{qoptions}
\opt{A}{The polymerization of chloroprene gives natural rubber.}
\opt{B}{Teflon is prepared from tetrafluoroethene by heating it with persulphate catalyst at high pressures.}
\opt{C}{Glyptal polymer is condensation polymer}
\opt{D}{Ethene at 350-570 K temperature and 1000-2000 atm pressure in the presence of a peroxide initiator yields high density polythene.}
\end{qoptions}
\end{bankq}
\begin{bankq}[id=SII23-P1-CHEM-S3-Q06, source={Score-II - Test 23 (17 Mar 2024)}, paper=1, qno=6, type=numerical, subtopic={Examples}, answer={Bonus}, answer_src=solved, working={Printed key omits this question's answer and it was not solved; banked as Bonus.}, difficulty=3, diff_parts={K1 S2 I0 C0 F1}, diff_note={concepts: condensation polymers | steps: terylene nylon 6,6 bakelite novolac are condensation polymers; cellulose is a natural addition of glucose by condensation but not counted; PVC SBR PTFE neoprene are addition; count}]
How many of the following are prepared by condensation polymerization?

(i) Decron (ii) Cellulose (iii) Nylon-6,6 (iv) PVC (v) SBR (vi) PTFE

(vii) Bakelite (viii) Novolac (ix) Neoprene
\end{bankq}

\section*{Score Advanced papers}
\begin{bankq}[id=SA03-P1-CHEM-S3-Q05, source={Score Advanced - Test 03 (27 Jun 2021)}, paper=1, qno=5, type=matching, subtopic={Types}, answer={P$\to$2, Q$\to$3, R$\to$1}, answer_src=solved, working={Terylene, Tollen's reagent, Buna-N}, difficulty=3, diff_parts={K1 S2 I0 C0 F1}, diff_note={concepts: polymers and reagents | steps: terylene is the condensation copolymer; Tollens reagent separates aldehydes; Buna-N addition copolymer}]
\noindent{\footnotesize\setlength{\tabcolsep}{2pt}\begin{tabular}{|L{.6cm}|L{3.2cm}|L{.6cm}|L{2.4cm}|}\hline
 & List-I & & List-II\\\hline
P & Synthetic condensation copolymer & 1 & Buna-N\\\hline
Q & Distinguish between aldehyde \& ketones & 2 & Terylene\\\hline
R & Addition copolymer & 3 & Tollen's reagent\\\hline
\end{tabular}}\par
\end{bankq}
\begin{bankq}[id=SA03-P2-CHEM-S1-Q07, source={Score Advanced - Test 03 (27 Jun 2021)}, paper=2, qno=7, type=single, subtopic={Examples}, answer={A}, answer_src=solved, working={Neoprene}, difficulty=2, diff_parts={K1 S1 I1 C0 F0}, diff_note={concepts: organic reaction product | steps: identify the reagent; product of the reaction}]
\par
\bfig{146.7}{SA03_p2_chem_s1_q07_1.png}
\begin{qoptions}
\opt{A}{\bfigi{104.0}{SA03_p2_chem_s1_q07_2.png}}
\opt{B}{\bfigi{102.4}{SA03_p2_chem_s1_q07_3.png}}
\opt{C}{\bfigi{104.0}{SA03_p2_chem_s1_q07_4.png}}
\opt{D}{\bfigi{102.5}{SA03_p2_chem_s1_q07_5.png}}
\end{qoptions}
\end{bankq}
\begin{bankq}[id=SA03-P2-CHEM-S1-Q08, source={Score Advanced - Test 03 (27 Jun 2021)}, paper=2, qno=8, type=single, subtopic={Types}, answer={B}, answer_src=solved, working={styrene}, difficulty=1, diff_parts={K1 S1 I0 C0 F0}, diff_note={concepts: monomer identification | steps: read the repeating unit; styrene}]
Identify the monomer in the given polymeric structures :\par
\bfig{74.9}{SA03_p2_chem_s1_q08_1.png}
\begin{qoptions}
\opt{A}{Propene}
\opt{B}{Styrene}
\opt{C}{Vinylchloride}
\opt{D}{Ethene}
\end{qoptions}
\end{bankq}
\begin{bankq}[id=SA05-P1-CHEM-S1-Q15, source={Score Advanced - Test 05 (15 Jul 2021)}, paper=1, qno=15, type=single, subtopic={Examples}, answer={D}, answer_src=key, difficulty=2, diff_parts={K1 S1 I0 C1 F0}, diff_note={concepts: monomers of Terylene | steps: polyester}]
Terylene (Dacron) is a condensation polymer of :
\begin{qoptions}
\opt{A}{Formaldehyde and urea}
\opt{B}{Ethylene glycol and ethylene diisocyanate}
\opt{C}{Ethylene glycol and phthalic acid}
\opt{D}{Ethylene glycol and terephthalic acidl}
\end{qoptions}
\end{bankq}
\begin{bankq}[id=SA06-P2-CHEM-S1-Q01, source={Score Advanced - Test 06 (04 Sep 2021)}, paper=2, qno=1, type=multi, subtopic={Examples}, answer={A,C}, answer_src=key, difficulty=3, diff_parts={K1 S2 I0 C0 F1}, diff_note={concepts: carbohydrate and polymer statements | steps: starch storage; natural rubber structure; third; fourth}]
Correct statement :
\begin{qoptions}
\opt{A}{Starch is the plant storage of carbohydrate.}
\opt{B}{Natural rubber is the trans-polymer.}
\opt{C}{Nylon show positive Lassigne test with Na metal.}
\opt{D}{Bakelite is phenol-acetaldehyde resin.}
\end{qoptions}
\end{bankq}
\begin{bankq}[id=SA06-P2-CHEM-S1-Q04, source={Score Advanced - Test 06 (04 Sep 2021)}, paper=2, qno=4, type=multi, subtopic={Types}, answer={C}, answer_src=key, difficulty=2, diff_parts={K1 S1 I0 C0 F1}, diff_note={concepts: biodegradable polymers | steps: PHBV and nylon-2-nylon-6 are biodegradable}]
Which of the following polymer is/are Biodegradable?
\begin{qoptions}
\opt{A}{\bfigi{84.9}{SA06_p2_chem_s1_q04_A.png}}
\opt{B}{\bfigi{133.9}{SA06_p2_chem_s1_q04_B.png}}
\opt{C}{\bfigi{148.0}{SA06_p2_chem_s1_q04_C.png}}
\opt{D}{\bfigi{128.5}{SA06_p2_chem_s1_q04_D.png}}
\end{qoptions}
\end{bankq}
\begin{bankq}[id=SA10-P2-CHEM-S1-Q04, source={Score Advanced - Test 10 (06 May 2023)}, paper=2, qno=4, type=single, subtopic={Types}, answer={D}, answer_src=key, difficulty=1, diff_parts={K1 S1 I0 C0 F0}, diff_note={concepts: monomers of PHBV | steps: 3-hydroxybutanoic acid and 3-hydroxypentanoic acid}]
Monomer units of Biodegradable polymer PHBV are :
\begin{qoptions}
\opt{A}{\bfigi{135.7}{SA10_p2_chem_s1_q04_A.png}}
\opt{B}{\bfigi{134.9}{SA10_p2_chem_s1_q04_B.png}}
\opt{C}{\bfigi{129.4}{SA10_p2_chem_s1_q04_C.png}}
\opt{D}{\bfigi{134.9}{SA10_p2_chem_s1_q04_D.png}}
\end{qoptions}
\end{bankq}
\begin{bankq}[id=SA16-P1-CHEM-S2-Q08, source={Score Advanced - Test 16 (20 May 2023)}, paper=1, qno=8, type=multi, subtopic={Mechanisms}, answer={A,B}, answer_src=key, difficulty=4, diff_parts={K1 S2 I1 C0 F1}, diff_note={concepts: monomers favourable for anionic polymerisation | steps: anionic polymerisation needs electron withdrawing groups: ester nitrile; vinyl acetate not; check}]
For which of the following monomers is anionic polymerization a favourable mechanism to form polymer ?
\begin{qoptions}
\opt{A}{CH$_2$ = CH $-$ COOCH$_3$}
\opt{B}{CH$_2$ = CH $-$ CN}
\opt{C}{CH$_2$ = CH $-$ OCOCH$_3$}
\opt{D}{CH$_2$ = C(CH$_3$)$_2$}
\end{qoptions}
\end{bankq}
\begin{bankq}[id=SA22-P1-CHEM-S3-Q04, source={Score Advanced - Test 22 (10 Apr 2024)}, paper=1, qno=4, type=numerical, subtopic={Types}, answer={4}, answer_src=key, difficulty=4, diff_parts={K1 S2 I0 C1 F1}, diff_note={concepts: count condensation polymers | steps: nylon 66 bakelite terylene silicones are condensation polymers; Teflon PVC Buna-S neoprene HDPE are addition; count}]
How many of the following is/are condensation polymer? Teflon, Nylon 66, PVC, Buna-S, Bakelite, Terylene, Neoprene, Silicones, HDPE.
\end{bankq}
\begin{bankq}[id=SA23-P2-CHEM-S1-Q04, source={Score Advanced - Test 23 (14 Apr 2024)}, paper=2, qno=4, type=multi, subtopic={Examples}, answer={A,D}, answer_src=key, difficulty=5, diff_parts={K1 S2 I1 C1 F1}, diff_note={concepts: polymers and monomers match | steps: natural rubber isoprene; Dacron terephthalic acid + ethylene glycol; nylon 6 caprolactam; check}]
Some name of polymers and corresponding monomers are given below. Identify the correctly matched option(s)?
\begin{qoptions}
\opt{A}{Natural rubber: \bfigi{56.8}{SA23_p2_chem_s1_q04_1.png}}
\opt{B}{Dacron: \bfigi{78.1}{SA23_p2_chem_s1_q04_2.png}}
\opt{C}{Nylon-6: \bfigi{35.0}{SA23_p2_chem_s1_q04_3.png}}
\opt{D}{Teflon: \bfigi{33.5}{SA23_p2_chem_s1_q04_4.png}}
\end{qoptions}
\end{bankq}
\begin{bankq}[id=SA28-P2-CHEM-S1-Q01, source={Score Advanced - Test 28 (27 Apr 2024)}, paper=2, qno=1, type=multi, subtopic={Examples}, answer={A,B,C,D}, answer_src=key, difficulty=6, diff_parts={K1 S3 I1 C1 F1}, diff_note={concepts: reaction scheme with step growth polymer and solubility | steps: identify compounds Q R from the scheme; Q has acid and phenol groups so dissolves in NaOH; solubility in NaHCO3 needs a strong acid; R formed by step growth polymerisation; check statements}]
\bfig{284.1}{SA28_p2_chem_s1_q01_1.png}

Choose correct option(s) ?
\begin{qoptions}
\opt{A}{(Q) is soluble in aq. NaHCO$_3$ as well as in aq. NaOH}
\opt{B}{Formation of (R) is step growth polymerisation}
\opt{C}{(P) will give yellow crystal with NaOI}
\opt{D}{(R) is polymer terylene}
\end{qoptions}
\end{bankq}
\begin{bankpassage}[id=SA36-P2-CHEM-S3-G3-4, source={Score Advanced - Test 36 (19 May 2024)}, questions={3-4}]
Given below is synthesis of nylon-66. Assume 100 \% yield in every step and none of the organic compounds involved are in excess.\par
\bfig{243.2}{SA36_p2_chem_s3_para3_4_1.png}
\end{bankpassage}
\begin{bankq}[id=SA36-P2-CHEM-S3-Q03, source={Score Advanced - Test 36 (19 May 2024)}, paper=2, qno=3, type=numerical, subtopic={Examples}, answer={130}, answer_src=key, difficulty=6, diff_parts={K2 S2 I1 C1 F1}, diff_note={concepts: moles or mass x in the synthesis, nylon 66 | steps: hexamethylenediamine and adipic acid condensation releases water; stoichiometry from the molar masses and the given mass; 100 percent yield at each step; evaluate x}, group=SA36-P2-CHEM-S3-G3-4]
The value of x is \underline{\hspace{1.5cm}}.
\end{bankq}
\begin{bankq}[id=SA36-P2-CHEM-S3-Q04, source={Score Advanced - Test 36 (19 May 2024)}, paper=2, qno=4, type=numerical, subtopic={Examples}, answer={365}, answer_src=key, difficulty=6, diff_parts={K2 S2 I1 C1 F1}, diff_note={concepts: value y, the nylon 66 synthesis | steps: apply stoichiometry from the previous step to the second quantity; moles of monomers; calculation of y}, group=SA36-P2-CHEM-S3-G3-4]
The value of y is \underline{\hspace{1.5cm}}.
\end{bankq}
\begin{bankpassage}[id=SA37-P2-CHEM-S3-G3-4, source={Score Advanced - Test 37 (19 May 2024)}, questions={3-4}]
Given below is synthesis of nylon-66. Assume 100 \% yield in every step and none of the organic compounds involved are in excess.\par
\bfig{243.2}{SA37_p2_chem_s3_para3_4_1.png}
\end{bankpassage}
\begin{bankq}[id=SA37-P2-CHEM-S3-Q03, source={Score Advanced - Test 37 (19 May 2024)}, paper=2, qno=3, type=numerical, subtopic={Examples}, answer={130}, answer_src=key, difficulty=5, diff_parts={K1 S2 I1 C1 F1}, diff_note={concepts: moles or mass x in the synthesis  nylon 66 | steps: hexamethylenediamine and adipic acid condensation releases water; stoichiometry from the molar masses and the given mass; 100 percent yield at each step; evaluate x}, group=SA37-P2-CHEM-S3-G3-4]
The value of x is \underline{\hspace{1.5cm}}.
\end{bankq}
\begin{bankq}[id=SA37-P2-CHEM-S3-Q04, source={Score Advanced - Test 37 (19 May 2024)}, paper=2, qno=4, type=numerical, subtopic={Examples}, answer={365}, answer_src=key, difficulty=5, diff_parts={K1 S2 I1 C1 F1}, diff_note={concepts: value y  the nylon 66 synthesis | steps: apply stoichiometry from the previous step to the second quantity; moles of monomers; calculation of y}, group=SA37-P2-CHEM-S3-G3-4]
The value of y is \underline{\hspace{1.5cm}}.
\end{bankq}
\begin{bankpassage}[id=SA38-P2-CHEM-S3-G1-2, source={Score Advanced - Test 38 (13 Apr 2025)}, questions={1-2}]
An organic compound \textbf P with molecular formula C$_9$H$_{18}$O$_2$ decolorizes bromine water and also shows positive iodoform test. \textbf P on ozonolysis followed by treatment with H$_2$O$_2$ gives \textbf Q and \textbf R. While compound \textbf Q shows positive iodoform test, compound \textbf R does not give positive iodoform test. \textbf Q and \textbf R on oxidation with pyridinium chlorochromate (PCC) followed by heating give \textbf S and \textbf T, respectively. Both \textbf S and \textbf T show positive iodoform test.\par
Complete copolymerization of 510 moles of \textbf Q and 510 moles of \textbf R gives one mole of a single acyclic copolymer \textbf U.\par
[Given, atomic mass : H = 1, C =12, O = 16]
\end{bankpassage}
\begin{bankq}[id=SA38-P2-CHEM-S3-Q01, source={Score Advanced - Test 38 (13 Apr 2025)}, paper=2, qno=1, type=numerical, subtopic={Examples}, answer={5}, answer_src=key, difficulty=6, diff_parts={K2 S2 I1 C1 F1}, diff_note={concepts: number of sp3 carbons in S, T from ozonolysis products | steps: P C9H18O2 decolourises bromine and gives iodoform: alkene with a CH3CH(OH) group; ozonolysis with H2O2 gives acids Q and R; further steps give S and T; count sp3 carbons in each}, group=SA38-P2-CHEM-S3-G1-2]
Sum of number of sp$^3$ carbon atoms in \textbf S and \textbf T is \underline{\hspace{1.5cm}}.
\end{bankq}
\begin{bankq}[id=SA38-P2-CHEM-S3-Q02, source={Score Advanced - Test 38 (13 Apr 2025)}, paper=2, qno=2, type=numerical, subtopic={Examples}, answer={5271}, answer_src=key, difficulty=6, diff_parts={K2 S2 I1 C1 F1}, diff_note={concepts: molecular weight, the copolymer divided by 18 | steps: identify the monomers from compounds Q and R; copolymer repeat unit mass from condensation or addition; U/18}, group=SA38-P2-CHEM-S3-G1-2]
The value of $\dfrac{\mathrm U}{18}$ is \underline{\hspace{1.5cm}}\par
where \textbf U is molecular weight of copolymer.
\end{bankq}
\begin{bankpassage}[id=SA39-P2-CHEM-S3-G1-2, source={Score Advanced - Test 39 (20 Apr 2025)}, questions={1-2}]
An acid A which is an important constituent of vineagar, on reaction with red P and Br$_2$ gives a monobromoderivative (B) which on reaction with NH$_3$ gives white solid (C). However ethanal on reaction with a mixture of ammonium chloride and sodium cyanide undergoes Strecker synthesis (synthesis of amino acids by the reaction of an aldehyde with cyanide in presence of ammonia) to give a product which on acidic hydrolysis gives another high melting solid (D). Let x and y are the total number of optical isomers of (C) and (D) respectively while z is the total number of dipeptides (excluding stereoisomers) that (C) and (D) can form.
\end{bankpassage}
\begin{bankq}[id=SA39-P2-CHEM-S3-Q01, source={Score Advanced - Test 39 (20 Apr 2025)}, paper=2, qno=1, type=numerical, subtopic={Examples}, answer={Bonus}, answer_src=solved, working={Printed key for this paper is incomplete and the figure/data needed to solve it was not recoverable; banked as Bonus.}, difficulty=6, diff_parts={K2 S2 I1 C1 F1}, diff_note={concepts: value, x + y in an organic scheme | steps: determine the structures along the sequence; count the required atoms or groups x and y; add}, group=SA39-P2-CHEM-S3-G1-2]
The value of x + y is \underline{\hspace{1.5cm}}.
\end{bankq}
\begin{bankq}[id=SA39-P2-CHEM-S3-Q02, source={Score Advanced - Test 39 (20 Apr 2025)}, paper=2, qno=2, type=numerical, subtopic={Examples}, answer={Bonus}, answer_src=solved, working={Printed key for this paper is incomplete and the figure/data needed to solve it was not recoverable; banked as Bonus.}, difficulty=5, diff_parts={K2 S1 I1 C1 F1}, diff_note={concepts: value, z in the same scheme | steps: from the structures obtained count the requested feature z}, group=SA39-P2-CHEM-S3-G1-2]
The value of z is \underline{\hspace{1.5cm}}.
\end{bankq}
\begin{bankq}[id=SA40-P2-CHEM-S3-Q07, source={Score Advanced - Test 40 (27 Apr 2025)}, paper=2, qno=7, type=matching, subtopic={Examples}, answer={Bonus}, answer_src=solved, working={Figure/data needed was not recoverable and no key given; banked as Bonus.}, difficulty=4, diff_parts={K1 S1 I1 C1 F1}, diff_note={concepts: match organic structures with their properties | steps: apply the property to each drawn structure; match}]
Match the Lists.\par
{\footnotesize\setlength{\tabcolsep}{1.5pt}\noindent\begin{tabular}{|c|c|c|c|}
\hline
 & \multicolumn{1}{c|}{List-I} & & \multicolumn{1}{c|}{List-II} \\ \hline
(P) & \bfigi{93.9}{SA40_p2_chem_s3_q07_1.png} & (1) & \bfigi{37.2}{SA40_p2_chem_s3_q07_5.png} \\ \hline
(Q) & \bfigi{112.1}{SA40_p2_chem_s3_q07_2.png} & (2) & \bfigi{55.4}{SA40_p2_chem_s3_q07_6.png} \\ \hline
(R) & \bfigi{99.2}{SA40_p2_chem_s3_q07_3.png} & (3) & \bfigi{30.1}{SA40_p2_chem_s3_q07_7.png} \\ \hline
(S) & \bfigi{96.2}{SA40_p2_chem_s3_q07_4.png} & (4) & \bfigi{33.5}{SA40_p2_chem_s3_q07_8.png} \\ \hline
\end{tabular}}\par
\end{bankq}
\begin{bankq}[id=SA40-P2-CHEM-S3-Q08, source={Score Advanced - Test 40 (27 Apr 2025)}, paper=2, qno=8, type=matching, subtopic={Examples}, answer={P$\to$4; Q$\to$3; R$\to$2; S$\to$1}, answer_src=key, difficulty=6, diff_parts={K2 S2 I1 C1 F1}, diff_note={concepts: pH, phosphate buffer mixtures matched with expressions | steps: moles after neutralisation: Na2HPO4 + HCl equal moles gives H2PO4- salt only; Na3PO4 + HCl 1.5 mmol : 3 mmol HCl gives; classify each mixture as buffer of pKa1 or pKa2 or giving H3PO4; match}]
Match the Lists.\par
{\footnotesize\setlength{\tabcolsep}{1.5pt}\noindent\begin{tabular}{|c|L{3.8cm}|c|L{2.8cm}|}
\hline
 & \multicolumn{1}{>{\centering\arraybackslash}p{3.8cm}|}{List-I} & & \multicolumn{1}{>{\centering\arraybackslash}p{2.8cm}|}{List-II (pH)} \\ \hline
(P) & 10 ml of 0.1 M Na$_2$HPO$_4$ + 10 ml of 0.1 M HCl & (1) & pH = pKa$_1$ + log $\dfrac{[\mathrm{salt}]}{[\mathrm{acid}]}$ \\ \hline
(Q) & 10 ml of 0.1 M Na$_3$PO$_4$ + 15 ml 0.2 M HCl & (2) & pH = pKa$_2$ + log $\dfrac{[\mathrm{salt}]}{[\mathrm{acid}]}$ \\ \hline
(R) & 10 ml of 0.1 M H$_3$PO$_4$ + 15 ml 0.1 M NaOH & (3) & H$_3$PO$_4$ is the product \\ \hline
(S) & 15 ml of 0.1 M Na$_3$PO$_4$ + 40 ml of 0.1 M HCl & (4) & pH = $\dfrac12$(pKa$_2$ + pKa$_1$) \\ \hline
\end{tabular}}\par
\end{bankq}

\section*{Enthusiast Pre-Score papers}
\begin{bankq}[id=EPS01-P2-CHEM-S2-Q06, source={Enthusiast Pre-Score - Test 01 (04 Nov 2020)}, paper=2, qno=6, type=single, subtopic={Types}, answer={5.00}, answer_src=key, difficulty=2, diff_parts={K1 S2 I0 C0 F0}, diff_note={concepts: polymer classification | steps: condensation versus addition; copolymer or homopolymer; count}]
Count number of polymers which are condensation polymers as well as co-polymers :\par
(1) Nylon-6 \quad (2) Nylon-6, 6 \quad (3) Nylon-6, 10 \quad (4) Dacron\par
(5) Bakelite \quad (6) Melamine formaldehyde resin \quad (7) Polyisoprene\par
(8) Styrene-butadiene rubber (Buna-S) \quad (9) PAN \quad (10) Teflon
\end{bankq}

\section*{JEE Advanced (2021 Series) papers}
\begin{bankq}[id=JA11-P1-CHEM-S1-Q02, source={JEE Advanced (2021 Series) - Test 11 (20 May 2021)}, paper=1, qno=2, type=multi, subtopic={Types}, answer={A,B,C,D}, answer_src=key, difficulty=7, diff_parts={K3 S3 I0 C1 F1}, diff_note={concepts: polymer classification, dipole moment of azulene, kinetic isotope effect in sulphonation | steps: bakelite is a condensation polymer of phenol and formaldehyde not a polyamide; biodegradable polymers like PHBV are synthetic; azulene is a polar hydrocarbon because of charge transfer to the 5-ring; sulphonation shows an isotope effect so rates differ; all four are incorrect | C: four unrelated statements from different units}]
Which of the following statement(s) is/are incorrect ?
\begin{qoptions}
\opt{A}{Bakelite is a polyamide}
\opt{B}{All Biodegradable polymers are natural polymers}
\opt{C}{Azulene \bfigi{54.4}{JA11_p1_chem_s1_q02_1.png} has zero dipole moment.}
\opt{D}{Sulphonation of C$_6$H$_6$, C$_6$D$_6$ \& C$_6$T$_6$ by Electrophilic aromatic substitution mechanism proceed at same rate.}
\end{qoptions}
\end{bankq}
\begin{bankq}[id=JA12-P1-CHEM-S1-Q10, source={JEE Advanced (2021 Series) - Test 12 (25 May 2021)}, paper=1, qno=10, type=multi, subtopic={Examples}, answer={A,C}, answer_src=solved, working={Cyclopentanone oxime by Beckmann rearrangement gives the six membered lactam for nylon-5; diethyl adipate gives cyclopentanone by Dieckmann condensation and hydrolysis; cyclohexanone gives nylon-6.}, difficulty=5, diff_parts={K2 S2 I1 C0 F1}, diff_note={concepts: Beckmann rearrangement of an oxime, Dieckmann condensation | steps: cyclopentanone oxime rearranges to the six membered lactam which polymerises to nylon-5; cyclohexanone gives caprolactam and nylon-6; diethyl adipate gives cyclopentanone through Dieckmann cyclisation hydrolysis and decarboxylation and then follows the first route | I: diethyl adipate must first be turned into cyclopentanone before the oxime route}]
By which of the following reaction sequence can Nylon--5 be prepared.
\begin{qoptions}
\opt{A}{\bfigi{162.9}{JA12_p1_chem_s1_q10_A.png}}
\opt{B}{\bfigi{164.2}{JA12_p1_chem_s1_q10_B.png}}
\opt{C}{\bfigi{163.6}{JA12_p1_chem_s1_q10_C.png}}
\opt{D}{All of these}
\end{qoptions}
\end{bankq}
\begin{bankq}[id=JA13-P1-CHEM-S2-Q01, source={JEE Advanced (2021 Series) - Test 13 (31 May 2021)}, paper=1, qno=1, type=single, subtopic={Types}, answer={C}, answer_src=solved, working={Nylon-6 monomer caprolactam shows amide resonance stabilisation.}, difficulty=3, diff_parts={K1 S2 I0 C1 F0}, diff_note={concepts: monomers of polymers and the chemistry of the monomers | steps: identify monomers: polystyrene from styrene, PVC from vinyl chloride, natural rubber from isoprene, nylon-6 from caprolactam; amide resonance stabilises caprolactam; check each combination | C: four rows against two other columns}]
Which of the following combination is correct?
\begin{qoptions}
\opt{A}{(III) (ii) (R)}
\opt{B}{(II) (i) (P)}
\opt{C}{(IV) (i) (P)}
\opt{D}{(I) (iii) (S)}
\end{qoptions}
\end{bankq}
\begin{bankq}[id=JA13-P1-CHEM-S2-Q02, source={JEE Advanced (2021 Series) - Test 13 (31 May 2021)}, paper=1, qno=2, type=single, subtopic={Types}, answer={D}, answer_src=solved, working={PVC monomer is vinyl chloride (iii) and not (iv).}, difficulty=1, diff_parts={K1 S1 I0 C0 F0}, diff_note={concepts: monomer of PVC and addition properties | steps: PVC is made from vinyl chloride (iii) and not from isoprene; the other combinations match monomers and properties correctly}]
Which of the following combination is incorrect?
\begin{qoptions}
\opt{A}{(I) (ii) (P)}
\opt{B}{(I) (ii) (R)}
\opt{C}{(III) (iv) (P)}
\opt{D}{(II) (iv) (R)}
\end{qoptions}
\end{bankq}
\begin{bankq}[id=JA13-P1-CHEM-S2-Q03, source={JEE Advanced (2021 Series) - Test 13 (31 May 2021)}, paper=1, qno=3, type=single, subtopic={Types}, answer={B}, answer_src=solved, working={Vinyl chloride adds HBr by Markovnikov 1,2 addition at room temperature.}, difficulty=2, diff_parts={K1 S2 I0 C0 F0}, diff_note={concepts: Markovnikov addition of HBr to vinyl chloride | steps: vinyl chloride is the monomer of PVC; HBr adds with Markovnikov orientation by 1,2 addition at room temperature; other options use the wrong monomer}]
Which of the following combination is correct?
\begin{qoptions}
\opt{A}{(II) (ii) (P)}
\opt{B}{(II) (iii) (R)}
\opt{C}{(III) (iii) (P)}
\opt{D}{(IV) (ii) (P)}
\end{qoptions}

\begin{center}\small\setlength{\tabcolsep}{4pt}
\begin{tabular}{|p{0.8cm}|p{3.8cm}|p{0.9cm}|p{4.6cm}|p{0.9cm}|p{3.6cm}|}\hline
 & \textbf{Column-I} & & \textbf{Column-II} & & \textbf{Column-III}\\\hline
(I) & Peptization & (i) & Colloid preparation technique & (P) & Sky in blue\\\hline
(II) & Delta formation & (ii) & Colloid destabilization & (Q) & Gold number\\\hline
(III) & Tyndall effect & (iii) & Colloid stabilization & (R) & Coagulation\\\hline
(IV) & Protection power of colloids & (iv) & Optical property of colloid & (S) & Adsorption of common ion\\\hline
\end{tabular}
\end{center}
\end{bankq}

\section*{TOAS papers}
\begin{bankq}[id=TOAS15-P1-CHEM-S2-Q10, source={TOAS - Test 15 (22 Aug 2022)}, paper=1, qno=10, type=multi, subtopic={Examples}, answer={A,B}, answer_src=key, difficulty=2, diff_parts={K1 S1 I0 C0 F1}, diff_note={concepts: phenol formaldehyde polymer | steps: identify the phenol among the reactions that forms it}]
$\text{P}+\text{HCHO}\xrightarrow{\text{OH}^-/\text{ Polymerisation}}$ Bakelite\par
Find the reaction in which P form as product.
\begin{qoptions}
\opt{A}{\bfigi{77.6}{TOAS15_p1_chem_s2_q10_A.png}}
\opt{B}{\bfigi{110.0}{TOAS15_p1_chem_s2_q10_B.png}}
\opt{C}{\bfigi{131.3}{TOAS15_p1_chem_s2_q10_C.png}}
\opt{D}{None}
\end{qoptions}
\end{bankq}

\section*{TMAS papers}
\begin{bankq}[id=TMAS04-P2-CHEM-S1-Q04, source={TMAS - Test 04 (19 Nov 2023)}, paper=2, qno=4, type=multi, subtopic={Examples}, answer={A,B,C,D}, answer_src=key, difficulty=3, diff_parts={K1 S1 I0 C1 F1}, diff_note={concepts: polymers | steps: check each}]
Which of the following statements are correct?
\begin{qoptions}
\opt{A}{Decron is co polymer of ethyleneglycol terephthatic acid}
\opt{B}{Nylon-6 is homo polymer of caprolactum}
\opt{C}{Novalac is co polymer of phenol and formaldehyde}
\opt{D}{Buna-S is the co polymer of butadiene and styrene.}
\end{qoptions}
\end{bankq}
\begin{bankq}[id=TMAS04-P2-CHEM-S3-Q04, source={TMAS - Test 04 (19 Nov 2023)}, paper=2, qno=4, type=numerical, subtopic={Types}, answer={5}, answer_src=key, difficulty=3, diff_parts={K1 S1 I0 C1 F1}, diff_note={concepts: condensation polymers | steps: classify each; count}]
How many of the following are condensation polymer?

(i) Nylon-6,6 \quad (ii) Buna-N \quad (iii) Nylon-6 \quad (iv) Terylene

(v) Teflon \quad (vi) Polystyrene \quad (vii) Glyptal \quad (viii) Bakelite
\end{bankq}

\section*{Star Review (Other) papers}
\begin{bankq}[id=STAR03-P2-CHEM-S4-Q07, source={Star Review (Other) - Test 03 (04 Dec 2023)}, paper=2, qno=7, type=numerical, subtopic={Types}, answer={5}, answer_src=key, difficulty=8, diff_parts={K1 S4 I2 C1 F1}, diff_note={concepts: true or false statements from several topics | steps: polymers; solubility; sugars; antiseptics; acidity; HVZ reaction; rates; isomerism; basicity; count}]
How many of the following statements are True ?\par
(a) Dacron is condensed copolymer\par
(b) Butan-1-ol is more soluble in water than propan-1-ol\par
(c) Lactose is reducing sugar but it gives negative Barfoed test\par
(d) $\mathrm{CHI_3}$ has antiseptic properties while $\mathrm{CHCl_3}$ has aneasthetic properties\par
(e) p-nitrophenol is more acidic than o-nitrophenol\par
(f) Carboxylic acid having $\alpha$-H shows HVZ reaction\par
(g) $\mathrm{CH_3{-}Br}$ has higher rate of $\mathrm{S_N2}$ than \bfigi{56.1}{STAR03_p2_chem_s4_q07_1.png}\par
(h) $1,1$-dichlorobut-$2$-ene can't show geometrical isomerism\par
(i) Secondary alkyl amine is less basic than tertiary alkyl amine in aqueous medium
\end{bankq}

\section*{Nurture (General) papers}
\begin{bankq}[id=NUR01-P1-CHEM-S1-Q02, source={Nurture (General) - Test 01 (12 Jun 2021)}, paper=1, qno=2, type=multi, subtopic={Types}, answer={A,B,C}, answer_src=solved, working={Bakelite is a condensation polymer of phenol-formaldehyde; PHBV is synthetic and biodegradable; azulene has a dipole; sulphonation shows no isotope effect}, difficulty=3, diff_parts={K1 S2 I0 C0 F1}, diff_note={concepts: polymer statements | steps: bakelite is a condensation polymer; biodegradable not all natural; third statement}]
Which of the following statement(s) is/are incorrect ?
\begin{qoptions}
\opt{A}{Bakelite is a polyamide}
\opt{B}{All Biodegradable polymers are natural polymers}
\opt{C}{Azulene \bfigi{68.6}{NUR01_p1_chem_s1_q02_1.png} has zero dipole moment.}
\opt{D}{Sulphonation of $\mathrm{C_6H_6}$, $\mathrm{C_6D_6}$ \& $\mathrm{C_6T_6}$ by Electrophilic aromatic substitution mechanism proceed at same rate.}
\end{qoptions}
\end{bankq}
\end{document}
